\par\medskip
\noindent\textbf{Teorema 5.3.}
Sea $X\longrightarrow Y$ un morfismo esencialmente de presentación
finita. Entonces:

\smallskip
\noindent\textbf{(1)} Para todo cambio de base
$Y'\longrightarrow Y$, la aplicación canónica
\[
J_{X'}^p(X'|Y')\xrightarrow{\ \sim\ }
J_X^p(X|Y)\times_YY'
\]
es un isomorfismo para todo $p$, donde $X'=X\times_YY'$.

\smallskip
\noindent\textbf{(2)} Para todo $x\in X$, existe un isomorfismo
canónico
\[
J_{X_x}^p(X_x|Y)\xrightarrow{\ \sim\ }J_X^p(X|Y)_x.
\]

\smallskip
\noindent\textbf{(3)} Hay inclusiones
\[
J_X^0(X|Y)\supset J_X^1(X|Y)\supset J_X^2(X|Y)\supset\cdots,
\]
y
\[
x\notin J_X^p(X|Y)
\qquad\text{para todo }p\geq
\dim_{\kappa(x)}\Omega_{X|Y}(x).
\]

\smallskip
\noindent\textbf{(4)} Un punto $x$ no pertenece a
$J_X^p(X|Y)$ si y solo si posee un entorno abierto $U$ que admite una
inmersión cerrada $i:U\hookrightarrow U'$ en un $Y$-esquema $U'$
tal que $U'$ es liso sobre $Y$ en $x'=i(x)$ y
\[
\dim_{x'}U'_{f(x)}=p.
\]

\smallskip
\noindent\textbf{(5)} Para las proyecciones
$\pi_i:X_1\times_YX_2\longrightarrow X_i$,
\[
\underline I_{X_1\times_YX_2}^r(X_1\times_YX_2|Y)
=\sum_{p+q=r}
 \pi_1^*\underline I_{X_1}^p(X_1|Y)\,
 \pi_2^*\underline I_{X_2}^q(X_2|Y).
\]

\smallskip
\noindent\textit{Demostración.}
La afirmación (3) se deduce de 5.1(b). La afirmación (1) se deduce de
5.1(a) y del isomorfismo canónico
\[
P^*\Omega_{X|Y}\xrightarrow{\ \sim\ }\Omega_{X'|Y'},
\qquad P:X\times_YY'\longrightarrow X.
\]
Asimismo, (2) se deduce de 5.1(a) y de
\[
\mathcal O_{X,x}\otimes_{\mathcal O_X}\Omega_{X|Y}
\xrightarrow{\ \sim\ }\Omega_{X_x|Y}.
\]
Por último, (5) se deduce de 5.1(d) y de
\[
\pi_1^*\Omega_{X_1|Y}\oplus\pi_2^*\Omega_{X_2|Y}
\xrightarrow{\ \sim\ }
\Omega_{X_1\times_YX_2|Y}.
\]

Para (4), supongamos primero que existe un diagrama
\[
\begin{tikzcd}[column sep=huge,row sep=large]
U \arrow[r,hook,"i"] \arrow[d,"f"']
  & U' \arrow[dl] \\
Y & {}
\end{tikzcd}
\]
donde $U'$ es liso sobre $Y$ en $x'=i(x)$ y
$\dim_{x'}U'_{f(x)}=p$. Entonces
\[
\dim_{\kappa(x')}\Omega_{U'|Y}(x')=p,
\]
y la sucesión
\[
i^*\Omega_{U'|Y}\longrightarrow\Omega_{U|Y}\longrightarrow0
\]
es exacta. Por tanto,
$\dim_{\kappa(x)}\Omega_{U|Y}(x)\leq p$, así que
$x\notin J_U^p(U|Y)$ por (3). Puesto que
\[
J_U^p(U|Y)=J_X^p(X|Y)|_U,
\]
tenemos $x\notin J_X^p(X|Y)$.

%%%%%%%%%%  scan idx 128  |  folio 117  %%%%%%%%%%

Recíprocamente, supongamos que $x\notin J_X^p(X|Y)$. La cuestión es
local, así que escribamos $X=\operatorname{Spec}(B)$ y
$Y=\operatorname{Spec}(A)$. Elijamos una $A$-presentación
\[
0\longrightarrow J\longrightarrow P\longrightarrow B
\longrightarrow0,
\qquad P=A[x_1,\ldots,x_n].
\]
Existe una sucesión exacta
\[
J/J^2\xrightarrow{d}B\otimes_P\Omega_{P|A}
\longrightarrow\Omega_{B|A}\longrightarrow0,
\]
y $B\otimes_P\Omega_{P|A}$ es libre de rango $n$ sobre $B$. Así,
\[
I_B^p(\Omega_{B|A})
=\operatorname{Im}\left(
 \bigwedge^{n-p}(J/J^2)\otimes_B
 (B\otimes_P\bigwedge^{n-p}\Omega_{P|A})^*
 \longrightarrow B\right).
\]
En consecuencia, $x\notin J_X^p(X|Y)$ precisamente cuando existen
$f_1,\ldots,f_{n-p}\in J$ y una elección de $n-p$ entre las variables
$x_1,\ldots,x_n$ para los cuales el menor jacobiano correspondiente
\[
\det\left(\frac{\partial f_i}{\partial x_j}\right)(x)
\]
es no nulo. Equivalentemente,
\[
\operatorname{Spec}\bigl(P/(f_1,\ldots,f_{n-p})\bigr)
\longrightarrow\operatorname{Spec}(A)
\]
es liso en la imagen $x'$ de $x$. Tomando
\[
U'=\operatorname{Spec}\bigl(P/(f_1,\ldots,f_{n-p})\bigr)
\]
se obtiene localmente la inmersión cerrada requerida. \hfill$\square$
